% ID: SEL-D0015-BC-E05
% ESTADO: cotejada. Segunda pasada visual completa y borrador compilado revisados.
% Fuente: Selectividad/matematicas_aplicadas/matematicas_aplicadas_P_C.pdf
% SHA-256: 463bf63d4697a18a0b9efc9287e7b0a593c803bab98a3cfab0e260067a0dd02c
% Página: 2; bloque C; ejercicio original 5.

A 120 estudiantes se les ha recomendado la lectura de dos libros. Se sabe
que 46 de ellos han leído el primer libro recomendado, 34 el segundo y 16
estudiantes han leído ambos libros. Se elige un estudiante al azar.
\begin{enumerate}[label=\alph*)]
  \item \textbf{(0.6 puntos)} Calcule la probabilidad de que haya leído alguno
  de los dos libros.
  \item \textbf{(0.6 puntos)} Calcule la probabilidad de que no haya leído
  ninguno de los dos libros.
  \item \textbf{(0.6 puntos)} Calcule la probabilidad de que solamente haya
  leído el primer libro.
  \item \textbf{(0.7 puntos)} Calcule la probabilidad de que haya leído el
  primer libro, si se sabe que no ha leído el segundo.
\end{enumerate}
